\documentclass[11pt]{article}
\usepackage{amsmath,amssymb,amsthm,mathtools,hyperref}

\begin{document}
\section*{Prebiotic polymerisation in water}
\label{prob:283}
\label{rank:189}
\par\smallskip

\noindent\textbf{Domain.} Prebiotic polymer chemistry; aqueous cycling. \quad \textbf{Retained status/source context.} Open functional-polymer accumulation question. Polymerization under wet--dry cycles has been demonstrated; a scalar polymerization rate exceeding hydrolysis is an inadequate general criterion.

\paragraph{Definitions and assumptions.} Let $c_{\ell,s}(t)$ be concentrations of polymers of length $\ell$ and linkage/sequence class $s$, with monomer feeds, ligation, hydrolysis and environmental loss represented in a finite truncation $1\leq\ell\leq L_{\max}$. Environment $e(t)$ includes water activity, pH and temperature and is either constant or cyclic with period $P$. Declare $P>0$ and the time origin fixing the cycle phase; for a constant environment, declare the positive averaging-window length $P$ used in the displayed criterion. The criterion takes the lower limit of successive window averages, not a cumulative average or an instantaneous lower bound. If the concentrations are stochastic, declare their law and the success quantifier for that criterion before evaluation; pathwise, almost-sure, expectation and probability requirements are distinct. A functional class $\mathcal F$ requires an experimentally tested template or catalytic property; merely detecting a long heterogeneous chain is insufficient. Specify linkage topology and measurement recovery bias. For a concrete instance, declare the polymer chemistry, length unit, concentration units and disjoint species classes, together with the independently assayed rule determining membership in $\mathcal F$. Explain how heterogeneous classes are accounted for without counting nonfunctional material as functional inventory, and distinguish function established under recovery/assay conditions from function during the environmental cycle.

\paragraph{Problem statement.} Find a physically justified environment for which the cycle-averaged functional inventory remains above a threshold:
\[
\liminf_{n\to\infty}\frac1P\int_{nP}^{(n+1)P}
\sum_{(\ell,s)\in\mathcal F}c_{\ell,s}(t)\,dt\geq c_{\min}>0,
\]
under declared continuing feed and losses. Validate the chain-length and linkage distributions, rather than infer them from a single net reaction constant.

\paragraph{Additional binding assumptions and resolution criterion.} Specify the treatment of polymers exceeding $L_{\max}$ and bound truncation error; reflecting or deleting overflow can artificially maintain functional inventory. State the error quantity and how its bound controls the effect of truncation on the displayed asymptotic functional-inventory criterion, accounting for return flux and interactions with omitted species where present. Distinguish finite-cycle observations from model-dependent justification of the infinite-time claim. The functional set is independently assayed and the positive threshold refers to a fixed volume convention. Declare the tracked region and phases, the amount-to-concentration conversion and its meaning throughout dry intervals, including how adsorption, precipitation, transport and rehydration affect the counted inventory. Specify whether the observable counts recoverable material or material available in a declared phase; these observables must not be interchanged. Require bounded physically feasible inventories and a declared continuing-energy/feed budget. A long-time averaged inventory is distinct from autonomous hereditary propagation of those polymers.

\paragraph{Alternative formulations.} (i) A constant-water formulation asks whether activated feeds or mineral surfaces sustain functional polymers without drying. (ii) A wet--dry formulation permits water loss and rehydration, while counting thermal degradation and material redistribution. They are not equivalent environments. A kinetic proof in a measured finite network is separate from a geochemical argument that the required feed and cycles occurred on early Earth.

\paragraph{Obstructions and partial progress.} Recent cyclic-nucleotide work demonstrates nonenzymatic mixed-sequence oligomer production and makes a blanket absence claim incorrect. Length, fidelity, linkage homogeneity and usable concentration remain separate requirements for heredity. Polymerization can occur predominantly during dry phases even though products are later aqueous; calling that spontaneous polymerization in bulk water would misdescribe the mechanism. The historical framing and protocell review do not verify all these rates. This problem concerns inventory and functionality, not yet self-sustained replication.

\paragraph{Sources (supplied provenance; not all rechecked).} \begin{itemize}
\item Malaterre, Jeancolas and Nghe, The Origin of Life: What Is the Question? (2022); conceptual framing, not problem-specific experimental verification. Metadata and indexed abstract were checked; direct page access was intermittent. \url{https://pubmed.ncbi.nlm.nih.gov/35594335/}
\item Lopez and Fiore, Investigating Prebiotic Protocells... (2019); review background; DOI fetch failed, abstract verified via PubMed PMID 31181679. \url{https://doi.org/10.3390/life9020049}
\item High-Yield Prebiotic Polymerization of 2-prime,3-prime-Cyclic Nucleotides under Wet-Dry Cycling (2025). \url{https://pubmed.ncbi.nlm.nih.gov/41019123/} (Primary oligomerization study; long functional replicators do not follow from yield alone.)
\end{itemize}


\section*{Binding catalogue conventions}
This edition revises both supplied catalogues and preserves original problem numbers 1--330. It adds one globally unique label \verb|\label{rank:x}| per problem, where smaller $x$ means greater provisional difficulty. A rank label is an anchor at that problem's section: \verb|\ref{rank:x}| returns the section number, not the literal difficulty number. The visible \texttt{rank:x} text and the index give the difficulty number itself.

The entries are \emph{parameterized mathematical research schemes}, including empirical prediction, causal identification, historical reconstruction, theory and engineering. They are not 330 certified unsolved mathematical conjectures. A rigorous instance fixes the objects and restrictions below; changing those objects changes the question. Established identities, finite calculations, non-identification constructions and solved special cases are not being posed anew as open. Some residual tasks concern informative bounds, biological realization or external validation rather than a new theorem.

Every main definition and problem statement received a mathematical editorial review. Each entry now has an additional binding-conditions paragraph; direct corrections are separately recorded in the accompanying audit. The alternative formulations inherit the applicable conventions, but are different targets unless equivalence is proved. Where an old broad phrase conflicts with an explicit restriction below it, the explicit restriction controls.

\textbf{Literature boundary.} Thirteen selective primary-literature checks cover eighteen entries in this revision; their access levels and scope are recorded explicitly. The supplied bibliography and its earlier source-access descriptions are retained as provenance, not represented as newly verified. Neither this review nor a source's description proves the absence of later solutions. The date of this edition is 13 October 2026; there is no claim of exhaustive current-status certification for all 330 entries.

\subsection*{Binding conventions for a concrete instance}
\paragraph{Objects, quantifiers and non-vacuity.}
An instance declares measurable state, input, output and parameter spaces; a fixed candidate family; admissible initial laws, interventions and environments; a finite horizon or specified asymptotic regime; observation and sampling laws; and performance metrics, positive normalization scales and positive tolerances. These are chosen before fitting or testing, not adjusted after seeing failures. Shared parameters must genuinely be shared across the declared validation domain. Finite dimension alone does not restrict arbitrary real-number encodings: specify a finite library, bounded description length, or regularity/complexity bounds. Exclude constant or no-signal solutions through a fixed informative-output or baseline-improvement requirement where appropriate. A universally quantified result cannot be established by choosing one trivial favorable instance.

\paragraph{Probability laws and data.}
Every expectation or probability refers to the specified stochastic law, including initial uncertainty, process noise, observation noise and sampling as appropriate. A deterministic equation is not a probability model without a declared random input or initial law. ``Observed'' targets in population-level formulas mean population observables $g_j^*$, not noisy estimates asserted to be exact. Let $\widetilde g_j$ be independent validation estimates and suppose
\[
\mathbb{P}\{\lvert\widetilde g_j-g_j^*\rvert\le r_j\text{ for every }j\}\ge1-\alpha.
\]
Then an uncertainty-aware sufficient certification rule for positive target tolerances $\varepsilon_j$ is
\[
\max_j\frac{|\widehat g_j-\widetilde g_j|+r_j}{\varepsilon_j}\le1.
\]
On that simultaneous confidence event it implies every population error is at most its tolerance, by the triangle inequality. Vector/norm versions use corresponding valid simultaneous bounds. Training uncertainty must also be included when the claim averages over repeated training data. Evaluation on finitely many conditions does not establish a continuous-domain supremum without covering and regularity bounds.

\paragraph{Dynamics and well-posedness.}
For ODEs and SDEs give coefficient domains, initial laws and existence/uniqueness conditions in the intended solution sense, or keep the entire admitted solution set. For stochastic jump models specify nonnegative off-diagonal rates, diagonal normalization, rate orientation and non-explosion. For PDEs specify spatial domains, boundary/initial conditions, conservation or source terms, and a weak or classical solution class. Moving-domain concentration equations include dilution/transport terms. A stationary law is used only for a time-homogeneous augmented model under existence assumptions; a stationary approximation to a driven system must be justified. Absorbing extinction states are retained. Survival-conditioned, quasi-stationary and unconditional laws are distinct and must not be substituted for one another.

\paragraph{Identification and causal meaning.}
Two parameter values are observationally equivalent when they induce the same full distribution under all allowed observation/intervention designs, not merely the same likelihood value for one sample. A target is identified only if constant on the equivalence class. Gauge symmetries, label switching and reference conventions are quotiented or fixed. Finite-data compatible sets use a declared confidence, likelihood or discrepancy rule; an empty set signals model rejection. Bayesian credible sets are not automatically confidence sets. A prior or regularizer does not create data identification. Bounds intended to resolve a question must be informative at a predeclared width or decision threshold.

A causal contrast uses well-defined intervention versions and a stated causal model. Random assignment, or explicit consistency, overlap and exchangeability assumptions, is required for causal interpretation from observations. Interference and time-varying confounding are represented when present. Fixing a mediator by intervention differs from conditioning on its observed value. Marginal potential-outcome laws do not determine individual cross-world effects or an arbitrary joint coupling. Ethically admissible designs, rather than unconstrained interventions, define the experimental menu.

\paragraph{Metrics, units, optimization and limits.}
Specify the sampling or reference measure for densities and integrals, the log base, and dimensionless arguments of logarithms. Total variation is $\sup_A|P(A)-Q(A)|$; Wasserstein distances require the declared ground metric and appropriate moment conditions; KL divergence uses its stated absolute-continuity convention. All normed readouts are scaled consistently, and error probabilities use $0<\alpha<1$. A ratio has a positive denominator or an explicit undefined/zero convention. Thermodynamic entropy production is nonnegative in expectation under the relevant hypotheses; individual stochastic realizations need not be. An infimum is not an attained minimum without compactness/coercivity and continuity or another existence argument. Fix the order of large-population, long-time and vanishing-mutation limits. Stability claims declare perturbation space and symmetries.

\subsection*{How the provisional difficulty order was assigned}
Four editorial components are scored from 1 (lower burden) to 5 (higher burden): (a) formal or inferential difficulty, (b) access to evidence and identification burden, (c) mechanistic/computational integration, and (d) independent validation and transfer. Their unweighted total orders the entries. For selected entries in the highest bands, an explicit editorial priority breaks equal totals; all other ties are broken by original problem number. Equal totals are \emph{not} evidence of reliably different true difficulties. Scores are judgments, not measurements, completion-time estimates or consensus rankings.

The object ranked is an informative, nondegenerate resolution of the \emph{scoped} scientific target, with the stated budget and validation domain. A sharp useful impossibility result can qualify; an already known identity, vacuous bound, arbitrary favorable model or lookup table cannot. Rank 330 is merely the least difficult relative assessment in this catalogue, not a claim of an easy biological solution. Broad historical titles do not automatically outrank precisely formulated mechanistic tasks.



\end{document}
